\begin{solapartat}
\figura[0.3]{sol_fig1.png}

\punts{0,2} Esquema

\punts{0,1} Per simetria, el camp total en la direcció vertical es cero.

Només hi ha camp elèctric en la direcció x:

Mòdul del camp elèctric en P:
\[ \left.\begin{aligned} &E_P = 2E\cos 60°\\ &E = k\frac{q}{l^2}\end{aligned}\right\}
E_P = \frac{2kq}{l^2}\cos 60° \qquad \text{\punts{0,2}} \]

% DUBTE: el càlcul dona 1686 N/C (1685,6); la pauta posa 1687 N/C.
\[ E_P = \frac{2\times 8,99\times 10^{9}\times 3,0\times 10^{-8}}{(0,40)^2}\cos 60° = 1687\ \mathrm{N/C} = 1,7\times 10^{3}\ \mathrm{N/C} \qquad \text{\punts{0,1}} \]

Direcció i sentit del camp elèctric: $\vec{E}_P = 1,7\times 10^{3}\,\vec{\imath}\ \mathrm{N/C}$ \punts{0,2}

\punts{0,2} $V_P = k\dfrac{q_+}{l} + k\dfrac{q_-}{l} = 0$

També es vàlida l'opció de col·locar la càrrega positiva a la dreta i la negativa a l'esquerra o qualsevol altra distribució que compleixi les condicions de l'enunciat.
\end{solapartat}

\begin{solapartat}
\punts{0,1} $U_i = k\dfrac{q_+q_-}{r_i}$

% DUBTE: el càlcul dona -2,02·10^-5 J (-2,0228·10^-5); la pauta posa -2,03·10^-5 J.
\punts{0,3} $U_i = 8,99\times 10^{9}\,\dfrac{3,0\times 10^{-8}(-3\times 10^{-8})}{0,40} = -2,03\times 10^{-5}\ \mathrm{J}$

Si es duplica la distància entre càrregues, l'energia potència és: $U_f = k\dfrac{q_+q_-}{r_f}$

on $r_f = 2\times 0,40\ \mathrm{m}$; $U_f = 8,99\times 10^{9}\,\dfrac{3,0\times 10^{-8}(-3\cdot 10^{-8})}{2\times 0,40} = -1,01\times 10^{-5}\ \mathrm{J}$ \punts{0,3}

$\Delta U = U_f - U_i \;\Rightarrow\; \Delta U = 1,01\times 10^{-5}\ \mathrm{J}$ \punts{0,3}. L'energia del sistema augmenta.
\end{solapartat}
